\def\topic{}
\ifdefined\MyHeaderIncluded
\else
\def\qwertyflag{true}

\ifdefined\MyHeaderIncluded

\else
  \newcommand{\foobar}{FooBar}
  \def\MyHeaderIncluded{true}
% if is till the end of the file

\documentclass[10pt,a4paper]{article}
\usepackage[utf8]{inputenc}
%\usepackage[german]{babel}
%\usepackage[cyrillic]{babel}
\usepackage[T1]{fontenc}
%\usepackage[T2A]{fontenc}
%\usepackage[english, russian]{babel}
\usepackage{amsmath}
\usepackage{amsfonts}
\usepackage{amssymb}
\usepackage{amsthm}
\usepackage{mathtools} %coloneqq

\usepackage{bussproofs} %

%\renewcommand{\thesubsection}{\arabic{subsection}}
%\newtheorem{thm}{Theorem}[section]
%\renewcommand{\thethm}{\arabic{subsection}.\arabic{thm}}

\usepackage{thmtools}
\usepackage{hyperref}
\hypersetup{
    colorlinks=true,
    linkcolor=blue,
    filecolor=magenta,      
    urlcolor=cyan,
    pdftitle={Overleaf Example},
    %pdfpagemode=FullScreen,
    }
%\usepackage{stix-type1}%\turnediota
\theoremstyle{definition}

\newtheorem{constructor}{Constructor}[section]
\newtheorem{iaxiom}{Axiom}
\newtheorem{axiom}{Axiom}
\newtheorem*{axiom*}{Axiom}
\newtheorem{definition}{Definition}[section]
\newtheorem*{definition*}{Definition}
\newtheorem{remark}{Remark}[section]
\newtheorem*{remark*}{Remark}
\newtheorem{recall}{Recall}[section]
\newtheorem*{recall*}{Recall}

\newtheorem{lemma}{Lemma}[section]
\newtheorem{theorem}{Theorem}[section]
\newtheorem*{theorem*}{Theorem}
\newtheorem{note}{Note}[section]
\usepackage{etoolbox}

%\newcommand*{\@Dashv}[2]{%
%  \reflectbox{$\m@th#1#2$}%
%}
\newcommand*{\separation}[3]{% 1)binder, 2)set, 3)condition
  \{#1\in #2: #3\}
}

% https://www.freecodecamp.org/news/simple-regex-tricks-for-beginners-3acb3fa257cb/
%  \\mathcal\{T\} |---> \topology
%  \\mathcal\{([^\}]*)\}   |--->  $1
\newcommand{\AND}{\ \ \&\ \ }
\newcommand{\any}{\mathrm{\ any}}
\newcommand{\borel}{\mathcal{B}}
\newcommand{\can}{\mathrm{can}}
\newcommand{\canonize}{\mathrm{canonize}}
\newcommand{\closure}{\mbox{cl}}
\newcommand{\complex}{\mathbb{C}}
\newcommand{\converges}{\longrightarrow}
\newcommand{\countable}{\mbox{\ countable}}
\newcommand{\ldefined}{\mbox{\ defined}}
\newcommand{\dom}{\mbox{dom}}
\newcommand{\Dom}{\mbox{Dom}}
\newcommand{\Done}{\noindent \textbf{Done.}\\}
\newcommand{\lexists}{\mbox{\ exists}}
\newcommand{\lforall}{\mbox{\ forall}}
\newcommand{\lunique}{\mbox{\ unique}}
\newcommand{\family}{\mathrm{\ family}}
\newcommand{\Fm}{\mathrm{Fm}}
\newcommand{\FuncSymb}{\mathrm{FuncSymb}}
\newcommand{\function}{\mathrm{\ function}}
\newcommand{\functional}{\mbox{\ functional}}
\newcommand{\IFF}{\ \Longleftrightarrow\ }
\renewcommand{\iff}{\ \longleftrightarrow\ }
\newcommand{\IMPLIES}{\ \Longrightarrow\ }
\renewcommand{\implies}{\ \longrightarrow\ }
\newcommand{\Inf}{\mathrm{Inf}}
\newcommand{\injective}{\mbox{\ injective}}
\newcommand{\injection}{\mbox{\ injection}}
\newcommand{\interior}{\mbox{int}}
\newcommand{\intersection}{\bigcap}
\newcommand{\isempty}{\mathrm{\ empty}}
\newcommand{\LB}{\mathrm{LB}}
\newcommand{\nghbset}{\mathcal{N}}
\newcommand{\nonempty}{\mathrm{\ nonempty}}
\newcommand{\Ord}{\mathrm{Ord}}
\newcommand{\power}{\mathcal{P}}
\newcommand{\pr}{\mathrm{pr}}
\newcommand{\fst}{\mathrm{fst}}
\newcommand{\snd}{\mathrm{snd}}
\newcommand{\PredSymb}{\mathrm{PredSymb}}
\newcommand{\Proof}{\noindent \textbf{Proof.}\\}
\newcommand{\ran}{\mathrm{ran}}
\newcommand{\Ran}{\mathrm{Ran}}
\newcommand{\real}{\mathbb{R}}
\newcommand{\relation}{\mathrm{\ relation}}
\newcommand{\set}{\mathrm{\ set}}
\newcommand{\singleton}{\mathrm{\ singleton}}
\newcommand{\semantics}{\mathrm{semantics}}
\newcommand{\semialgebra}{\mathcal{S}}
\newcommand{\Sup}{\mathrm{Sup}}
\newcommand{\surjective}{\mbox{\ surjective\ on\ }}
\newcommand{\Term}{\mathrm{Term}}
\newcommand{\Type}{\mathrm{Type}}
\newcommand{\topology}{\mathcal{T}}
\newcommand{\total}{\mbox{\ total\ on\ }}
\newcommand{\transitive}{\mathrm{\ transitive}}
\newcommand{\union}{\bigcup}
\newcommand{\lTrue}{\mathrm{True}}
\newcommand{\lFalse}{\mathrm{False}}
\newcommand{\lUncertain}{\mathrm{Uncertain}}
\newcommand{\lif}{\mathrm{if}}
\newcommand{\lthen}{\mathrm{then}}
\newcommand{\lelse}{\mathrm{else}}
\newcommand{\lexception}{\mathrm{exception}}
\newcommand{\lundef}{\mathrm{undef}}
\newcommand{\univocal}{\mathrm{\ univocal}}
\newcommand{\univocalat}{\mathrm{\ univocal\ at\ }}
\newcommand{\UB}{\mathrm{UB}}
\newcommand{\Un}{\mathcal{U}}
\newcommand{\Var}{\mathrm{Var}}

\newcommand{\notvarphi}{\phi}
\renewcommand{\phi}{\varphi}

\newcommand{\comment}[1]{\quad\mbox{(#1)}}
% local commands
%\preto\section{\undef\blipo} %etoolbox
\preto\subsection{\undef\blipo} %etoolbox
% https://tex.stackexchange.com/questions/105684/limit-macro-newcommand-scope-to-section
\usepackage{graphicx}
\usepackage{authblk}

\setlength\parindent{0pt} % whole doc of \noindent
\usepackage{soul}
%\usepackage{mathgl}

\fi % guard
\ifdefined\topic
\else
\def\topic{No Titel}
\fi
\title{\topic}
\author[1]{Georgy Dunaev}
\begin{document}
\maketitle
\fi
\section{Sethood and elementary class reasoning}
This section derives basic membership and inclusion facts from
comprehension, extensionality, and the definition of sethood. The notation
$X\set$ means that $X$ is a member of some class. In particular, $\V$
contains all sets, while every class is a subclass of $\V$; these are
statements about membership and inclusion respectively.

\iftrue
%The following theorem requires finding a fresh variable. But since it is always possible, the variant of it without the side contion may be added as and axiom, just to reduce the proofs.
%%%%%%%%%%%%%%%%%%%%%%%%%%%%%%%%%%%%%%%%% COMMENTED OUT
\begin{theorem}[example of class builder use]
\[
(\cc{y}{\psi}\nonempty)
\iff
(\exists (y\set),\ \psi)
\]
\end{theorem}
\begin{proof}
A class is nonempty exactly when it has a member. By comprehension,
$z\in\cc{y}{\psi}$ holds exactly when $z$ is a set and $\psi[z/y]$ holds.
Thus a member of the comprehension supplies the required witness, and any
such set witness is a member of the comprehension.
\end{proof}
\fi

\begin{theorem}[V is maximal]
Every class is a subclass of the class of all sets.
\[
(C\subseteq \V)
\]
\end{theorem}
\Proof
%\[\forall x\in C,\ x\in \V\]
%...
Claim: $(\forall x\in C,\ x\in \V)$.\\
Fix $x$. Assume $(x\in C)$. To show $(x\in \V)$.
\[
(x\in C)
\IMPLIES
(x\set)
\tag{by \myref{thm:in class then set}}
\]
\[
(x\set)
\LAND
\top
\IMPLIES
(x\in \V)
\tag{by \myref{def:universal class}}
\]
The claim is proven.
\[
(\forall x\in C,\ x\in \V)
\IMPLIES
(C\subseteq \V)
\tag{by \myref{def:subclass}}
\]
\Done

\begin{theorem}[sets are in V]
\[
(X\set)\iff(X\in \V)
\label{thm:sets are in V}
\]
\end{theorem}
\Proof
\noindent($\Leftarrow$) Assume $(X\in \V)$.
\[
(X\in \V)\IMPLIES(\exists W,\ X\in W).
\]
\[
(\exists W,\ X\in W)\IMPLIES(X\set).
\]
\noindent($\Rightarrow$) Assume $(X\set)$.
%\[
%(X\set)\IMPLIES(\exists W,\  X\in W))
%\]
\[
(X\set)\LAND\top[X/W]\IMPLIES(X\in \cc{x}{\top[x/W]})
\]
\[
(X\in \cc{x}{\top[x/W]}=\cc{x}{\top}=\V)\IMPLIES (X\in \V)
\]
%\[(a\in\{a\})\LAND(\{a\}=\{w\})\IMPLIES(a\in\{w\})\]
\Done
\begin{theorem}[antisymmetry of inclusion]
\[
(A \subseteq B)\land (B\subseteq A) \implies (A=B)
\]
\end{theorem}
\Proof
Claim: $(\forall x,\ x\in A \iff x\in B)$.\\
Fix x.\\
$(\Rightarrow)$\\
Assume $(x\in A)$. To show: $(x\in B)$\\
\[
(x\in A)
\LAND
(A\subseteq B)
\IMPLIES
(x\in B)
\tag{by \myref{thm:element of subclass}}
\]
$(\Leftarrow)$\\
Assume $(x\in B)$. To show: $(x\in A)$\\
\[
(x\in B)
\LAND
(B\subseteq A)
\IMPLIES
(x\in A)
\tag{by thm \myref{thm:element of subclass}}
\]
The claim has been proven.
\[
(\forall x,\ x\in A \iff x\in B)
\IMPLIES
(A=B)
\tag{by \myref{axm:extensionality}}
\]
\Done

\begin{theorem}[empty set is minimal]
For every class $A$,
\[
(\emptyclass\subseteq A)
\clabeltagd{thm}{empty class is subclass}
\]
\end{theorem}
\Proof
Claim:
\[
(\forall x\in \emptyclass,\  x\in A)
\]
Assume $(x\in\emptyclass)$.
\[
(x\in\emptyclass)
\IMPLIES
\bot
\]
\[
\bot
\IMPLIES
(x\in A)
\]
The claim is proven.
\[
(\forall x\in \emptyclass, x\in A)
\IMPLIES
(\emptyclass\subseteq A)
\tag{by \myref{def:emptyclass}}
\]
\Done

\begin{theorem}
\[
\neg(\exists w, w\in A)
\implies
(A=\emptyclass)
\]
\end{theorem}
\Proof
Assume $\neg\exists w,\  w\in A$. To show: $(A=\emptyclass)$.\\
$(\supseteq)$\\
\[
(A\supseteq\emptyclass)
\tag{by \myref{thm:empty class is subclass}}
\]
$(\subseteq)$\\
Claim: $(\forall x\in A,\ x\in\emptyclass)$\\
Fix $x$. Assume $(x\in A)$. To show: $\bot$\\
\[
(x\in A)
\IMPLIES
(\exists x,\ x\in A)
\]
\[
(\exists x,\ x\in A)
\LAND
(\neg\exists x,\ x\in A)
\IMPLIES
\bot
\]
\[
\bot
\IMPLIES
(x\in\emptyclass)
\]
The claim is proven.
\[
(\forall x\in A,\ x\in\emptyclass)
\IMPLIES
(A \subseteq \emptyclass)
\tag{by def of subclass \myref{def:subclass}}
\]
\[
(A \subseteq \emptyclass)
\LAND
(A \supseteq \emptyclass)
\IMPLIES
(A=\emptyclass)
\tag{by axm \myref{axm:extensionality}}
\]
%So by extensionality axiom we are done.\\
%\[
%\neg\exists w,\  w\in A
%\IMPLIES
%(A=\emptyclass)
%\]
\Done
\begin{theorem}[emptiness criterion]
\[
\neg\exists w,\  w\in A
\iff
(A\isempty)
\]
\end{theorem}
\Proof
By the preceding theorem, having no member implies equality with
$\emptyclass$. Conversely, equality with $\emptyclass$ implies having
no member, since the empty class has none. This is exactly the stated
criterion for $A\isempty$.
\Done
\begin{theorem}[nonemptiness criterion]
\[
\exists w,\  w\in A
\iff
(A\nonempty)
\]
\end{theorem}
\Proof
Negate the emptiness criterion and use classical double-negation
elimination. Thus $A\ne\emptyclass$ holds exactly when
$\exists w\,(w\in A)$.
\Done

\begin{theorem}[membership in a subclass]
Any element of a subclass is in the class.
\[
(x\in A)\land(A\subseteq B)\implies(x\in B)
\label{thm:element of subclass}
\]
\end{theorem}
\ProofNoPar
\[
(A\subseteq B)
\IMPLIES
(\forall x\in A,\ x\in B)
\tag{by def of subclass \myref{def:subclass}}
\]
\[
(x\in A)
\LAND
(\forall x\in A,\ x\in B)
\IMPLIES
(x\in B)
\tag{by restricted specification elimination}
\]
\Done

\begin{theorem}[is in class then is set]
For any terms $A$ and $B$,
%\[A\set \land B\in A \rightarrow B\set\]
\[
(B\in A)
\implies
(B\set)
\]
\label{thm:in class then set}
\end{theorem}
\Proof
Assume $(B\in A)$. To show: $(B\set)$. Choose a variable $Y$
not free in either term $A$ or $B$.
\[
(B\in A)
\IMPLIES
(\exists Y,\  B\in Y)
\tag{by right--exists--intro}.
\]
\[
(\exists Y,\  B\in Y)
\IMPLIES
(B\set)
\tag{introduced by set's definition}.
\]
\Done

\begin{theorem}[a class with exactly one member]
\[
(w\in A) \land (!a,\ a\in A) \implies (\{w\}=A)
\]
\end{theorem}
\Proof Assume $(w\in A)$ and $(!a,\ a\in A)$. Membership implies
that $w$ and every other member of $A$ are sets, so singleton introduction
is legitimate in the following argument.\\
\noindent($\subseteq$) Assume $(q\in \{w\})$.
\[
(q\in \{w\})\IMPLIES(q=w).
\]
\[
(w\in A)\LAND(q=w)\IMPLIES(q\in A).
\]
\noindent($\supseteq$) Assume $(a\in A)$.
\[
(!a,\ a\in A)\LAND(a\in A)\LAND(w\in A)\IMPLIES(a=w).
\]
\[
(a=w)\IMPLIES(\{a\}=\{w\})
\]
\[
(a\in\{a\})\LAND(\{a\}=\{w\})\IMPLIES(a\in\{w\})
\]
\Done

\ifdefined\qwertyflag
\undef\qwertyflag
\end{document}
\fi
